\documentclass[12pt]{article}
%packages
\usepackage{amsmath, amsthm, amssymb}
\usepackage{graphicx}
\usepackage{enumerate}
\usepackage[utf8]{inputenc}
% \usepackage[spaced=0]{diffcoeff}
\usepackage[italic=true]{derivative}
\usepackage{tcolorbox}
\usepackage{physics2}
\usephysicsmodule{ab, braket}
\usepackage{siunitx}
\usepackage{fancyhdr}
\usepackage[margin=1in]{geometry}
\usepackage{tikz}
\usepackage{pgfplots} % for the axis environment
% \usepackage{fullpage}

%environments
\newtheorem*{question}{Q}
\newtheorem*{prop}{Proposition}

\newenvironment{quest}[1]
{\vspace{0.3cm}\noindent\textbf{Q(#1)}\em}{\vspace{0.3cm}}

\newenvironment*{Remark}{\noindent\textbf{Remarks}\begin{enumerate}}{\end{enumerate}\vspace{1em}\noindent}


\usepackage{hyperref}
\hypersetup{
    colorlinks=true,
    linkcolor=blue,
    filecolor=magenta,      
    urlcolor=cyan,
    pdftitle={Overleaf Example},
    pdfpagemode=FullScreen,
}


%commands
\newcommand{\e}{\varepsilon}
\newcommand{\R}{\mathbb{R}}
\newcommand{\Z}{\mathbb{Z}}
\newcommand{\N}{\mathbb{N}}
\newcommand{\C}{\mathbb{C}}
\newcommand{\Q}{\mathbb{Q}}


\newcommand{\grad}{\vec \nabla}
\newcommand{\LL}{\mathcal{L}}
\newcommand{\vphi}{\varphi}
\newcommand{\p}{\partial}



\newcommand{\eq}{\Leftrightarrow}
\newcommand{\clm}{\textbf{Claim }}
\newcommand{\rmks}{\noindent\textbf{Remarks}}

\newcommand{\levi}{\mathfrak{E}}


\makeatletter
     \newcommand\vb{\@ifstar\boldsymbol\mathbf}
     \newcommand\va[1]{\@ifstar{\vec{#1}}{\vec{\mathrm{#1}}}}
     \newcommand\vu[1]{%
       \@ifstar{\hat{{#1}}}{\hat{{#1}}}}
    % \renewcommand*{\vec}[][default]{def}
\makeatother


\newcommand{\ev}[1]{\langle{#1}\rangle}
\newcommand{\vecp}[1]{\vec{#1}\:'}
%\renewcommand\qedsymbol{QED}


\newcommand{\HH}{\mathcal{H}}
\newcommand{\norm}[1]{||#1||}
\newcommand{\pv}{\mathcal{P}}
\newcommand{\ubra}[1]{\langle #1 \vert}
\newcommand{\uket}[1]{\vert #1 \rangle}
\newcommand{\ubraket}[2]{\langle #1 \vert #2 \rangle}
\newcommand{\uelem}[3]{\langle #1 \vert #2 \vert  #3 \rangle}

%%%%%%%%%%%%%%%%%%%%%%
% Set up header/footer
\pagestyle{fancy}
\fancyhead[LO,L]{Vishwanath Wimalasena}
\setlength{\headheight}{14.5pt}
\fancyfoot[LO,L]{}
\fancyfoot[CO,C]{\thepage}
\fancyfoot[RO,R]{}
\renewcommand{\headrulewidth}{0.5pt}
\renewcommand{\footrulewidth}{0.4pt}
\renewcommand{\headrulewidth}{0.5pt}
%%%%%%%%%%%%%%%%%%%%%%



\tcbuselibrary{theorems}
\tcbuselibrary{breakable}
\tcbuselibrary{skins}

%environments
\newtheorem{theorem}{Theorem}[section]
\newtheorem{proposition}{Proposition}[section]
\newtheorem{definition}{Definition}[section]
\newtheorem{lemma}{Lemma}[section]
\newtheorem{lem}{Lemma}
\newtheorem*{ex}{Exercise}

\theoremstyle{definition}
\newtheorem*{sol}{Sol}



% \newtcolorbox{ansbox}{tcbox raise base, leftrule=0pt, arc=0pt, rightrule=0pt,toprule=0mm,bottomrule=0mm, colback=green!8!white, colframe=green!20!black}

% \newtcolorbox{background}[1][]{enhanced,
%     colback=green!5!white,colframe=green!40!black,
%     colbacktitle=green!40!black,fonttitle=\bfseries,
%     frame hidden,
%     title=Background,
%     boxrule=0pt,
%     titlerule style=green!70!black, #1}

\tcolorboxenvironment{ex}{enhanced jigsaw, colback=green!8!white, colframe=green!40!black,breakable,before skip=10pt,after skip=10pt}

\tcolorboxenvironment{question}{% `proof' from `amsthm' 
    blanker,colback=green,breakable,left=5mm,
    before skip=10pt,after skip=10pt,
    borderline west={1mm}{0pt}{green!40!black}}

% \tcbset{one-sided-color/.style={
%         breakable,
%         enhanced,
%         boxrule=0pt,
%         frame hidden,
%         borderline ={1pt}{0pt}{#1},
%         colback=#1!10,
%         coltitle=#1!200,
%         after title={\ },
%         attach title  to upper,
%     }
% }

\newtcolorbox{q}{one-sided-color=green!60!black, title=\textbf{Q}}

\newtcolorbox{outline}[1]{one-sided-color=blue!70!black, title=\textbf{#1}}

% \newtcolorbox{ans}{breakable,
%         enhanced,
%         boxrule=0pt,
%         frame hidden,
%         borderline west={1pt}{0pt}{green!70!black},
%         % borderline south={1pt}{0pt}{green!70!black},
%         % after title={\ },
%         % attach title  to upper,
%         top=1pt,
%         % boxsep=1pt,
%         colback=white,
%         coltitle=green!70!black!200}


% \newtcolorbox{q2}{breakable,
%         enhanced,
%         boxrule=0pt,
%         frame hidden,
%         borderline west={1pt}{0pt}{green!70!black},
%         colback=green!70!black!10, after title={Q}}

% \newtcolorbox{exc}[1]{one-sided-color=green!70!black, title=\textbf{(#1)}}
\definecolor{bordercolor}{rgb}{0.13,0.57,0.07}
\definecolor{boxcolor}{rgb}{0.28,0.69,0.28}

\newtcolorbox{exc}[1]{
        breakable,
        enhanced,
        frame hidden,
        borderline ={1pt}{0pt}{bordercolor},
        colback=boxcolor!10,
        coltitle=bordercolor,
        after title={\ },
        attach title  to upper,
        title=\textbf{Exercise #1},
        }

\newtcolorbox{ans}{
    breakable,
    enhanced,
    boxrule=0pt,
    frame hidden,
    borderline west={1pt}{0pt}{bordercolor},
    top=0pt,
    colback=white,
    coltitle=bordercolor,
    }
% \newcommand{\exc}[1]{\begin{exc}#1\end{exc}}


\newtcbtheorem[number within=section]{cthm}{Theorem}%
{colback=green!8!white, colframe=green!40!black,fonttitle=\bfseries, fontupper=\itshape}{mainthm}

\newtcbtheorem[number within=section]{cexercise}{Exercise}%
{colback=green!8!white, colframe=green!40!black,fonttitle=\bfseries, fontupper=\itshape}{mainthm}

\newtcbtheorem[number within=section]{cclaim}{Claim}%
{colback=green!8!white,colframe=green!40!black,fonttitle=\bfseries, fontupper=\itshape}{claim}

\tcolorboxenvironment{sol}{breakable, enhanced, left=5mm,
    before skip=10pt, after skip=10pt, frame  hidden,colback=green!6!white,
    borderline west={1pt}{0pt}{green!60!black}}

\tcolorboxenvironment{lem}{
    enhanced jigsaw,colback=green!8!white, colframe=bordercolor, colback=white, breakable,before skip=10pt,after skip=10pt }

\tcolorboxenvironment{proof}{% `proof' from `amsthm' 
    blanker,breakable,left=5mm,
    before skip=10pt,after skip=10pt,
    borderline west={1pt}{0pt}{green!60!black}}

\fancyhead[CO,C]{\textbf{PHYS606}}
\fancyhead[RO,R]{PSET 2}

\begin{document}
\begin{exc}{1a}
\end{exc}
\begin{ans}
    The full hamiltonian for the system is given by 
    \[\HH = \HH_{\text{atom}} + \HH_{\text{int}} + \HH_{\text{field}} \]
    Let $\ketbra{i}{j} = \ket{i} \otimes \bra{j_k}$ be a basis for the system with $\ket{i}$ being eigenstates of the atom and $\ket{j_k}$ representing a field with $j$ photons in mode $k$. Setting $\ket{1}$ to be zero energy, the atomic hamiltonian is given by
    \begin{align*}
        \HH_{\text{atom}}&= \hbar\omega_{12}\ketbra{2}{2} + \hbar\omega_{13}\ketbra{3}{3}
    \end{align*}
    The radiation field has hamiltonian 
    \begin{align*}
        \HH_{\text{field}} &= \sum_{k, \lambda} \hbar\nu_{k,\lambda}\ab(a_{k\lambda}^\dagger a_{k\lambda} + \frac{1}{2})
    \end{align*}
    and the interaction hamiltonian is given by 
    \begin{align*}
        \HH_{\text{int}} &= - \hat d \cdot \hat E = -\sum_{\alpha, \alpha'} \ketbra{\alpha}{\alpha} \hat d \ketbra{\alpha'}{\alpha'} \cdot \sum_{k, \lambda} \sqrt{\frac{\hbar \nu_k}{2\e_0 V}} \left(\hat a_{k\lambda}\hat \e_{k\lambda}e^{i\vec k \cdot \vec r} + \hat a_{k\lambda}^\dagger \e_{k\lambda}^* e^{-i\vec k \cdot \vec r}\right)
    \end{align*}
    By assumption, $\uelem{2}{\hat d}{1} = \uelem{3}{\hat d}{1} = \mu \vu z$. Letting $g_k = \frac{\mu}{\hbar}\sqrt{\hbar \nu_k/2\e_0 V}(\hat \e_k \cdot \hat z) e^{i\vec k \cdot \vec r}$, we get  
    \begin{align*}
        \HH_{\text{int}} &= -\hbar\sum_{k, \lambda} g_k \left(\ketbra{2}{1} + \ketbra{3}{1}\right) \hat a_{k\lambda} + g_k^* \left(\ketbra{1}{2} + \ketbra{1}{3}\right) \hat a_{k\lambda}^\dagger
    \end{align*}
    To study spontaneous emission, we consider states $\ket{2, 0}, \ket{3, 0}$ (atom in excited state, no photons) and $\ket{1, 1_k}$ (atom in ground state after emitting photon in mode $k$). We evaluate the matrix elements as 
    \begin{align*}
        \uelem{2, 0}{\HH}{2, 0} &= \hbar \omega_{12}\\
        \uelem{3, 0}{\HH}{3, 0} &= \hbar \omega_{13}\\
        \uelem{1, 1_k}{\HH}{1, 1_k} &= \hbar \nu_k \\
        \uelem{2, 0}{\HH}{1, 1_k} &= -\hbar g_k \\
        \uelem{3, 0}{\HH}{1, 1_k} &= -\hbar g_k
    \end{align*}
    where we ignored the $\hbar\nu_k/2$ term in the field hamiltonian since it adds a constant to the total energy, irrelevant to dynamics. Thus, we get 
    \begin{align*}
        \HH &= \underbrace{\hbar\omega_{12}\ketbra{2}{2} + \hbar\omega_{13}\ketbra{3}{3} + \sum_{k}\hbar \nu_k\ketbra{k}{k}}_{:= \HH_0}  - \sum_{k} \ab(\hbar g_k\ketbra{2}{k} + \hbar g_k\ketbra{3}{k} + \text{h.c.})
    \end{align*}
    where we dropped the sum over $\lambda$ by absorbing it implicitly into the sum over $k$.  Let $\HH_1$ be the interaction part of the hamiltonian, i.e $\HH_1 = \HH - \HH_0$. We move to the interaction picture w.r.t $\HH_0$ by defining $\ket{\psi_I} = \hat U^{-1}\ket{\psi}$ where $\hat U = e^{-i\HH_0 t/\hbar}$. Then, 
    \begin{align*}
        i\hbar\odv{\ket{\psi_I}}{t} &= i\hbar \odv{\hat U^{-1}}{t} \ket{\psi} + i\hbar \hat{U}^{-1}\odv{\ket{\psi}}{t} \\
        &= - \HH_0 \hat U^{-1} \ket{\psi} + \hat U^{-1}(\HH_0 + \HH_1) \ket\psi \\
        &= \hat U^{-1} \HH_1 \ket{\psi}\\
        &= \hat U^{-1} \HH_1 \hat U \ket{\psi_I}
    \end{align*}
    where we used the fact that $\hat U^{-1}$ commutes with $\HH_0$ in the last step. Thus, our effective hamiltonian in the interaction picture is $\HH_{\text{eff}} = \hat U^{-1} \HH_1 \hat U$. We note that since $\HH_0$ is diagonal, we have 
    \begin{align*}
        \hat U^{-1}= e^{i\HH_0 t/\hbar} = e^{i\omega_{12}t}\ketbra{2}{2} + e^{i\omega_{13}t}\ketbra{3}{3} + \sum_k e^{i\nu_k t}\ketbra{k}{k}
    \end{align*}
    Hence, we get
    \begin{align*}
        \HH_{\text{eff}} &= -\hbar \sum_k g_k \ketbra{2}{k} e^{-i(\nu_k - \omega_{12})t} + g_k \ketbra{3}{k} e^{-i(\nu_k - \omega_{13})t}+ \text{h.c.}
    \end{align*}
\end{ans}

\begin{exc}{1b}
    
\end{exc}
\begin{ans}
    Apply the unitary operator $\hat U = e^{-i\sum_k (\nu_k - \omega_{12}) \ketbra{k}{k}}$
    We write $\ket{\psi_I(t)} = c_2(t) \ket 2 + c_3(t) \ket 3 + \sum_k c_k(t) \ket{k}$. Using the effective hamiltonian from part (a), we find 
    \begin{align*}
        i\hbar \odv{}{t}\ket{\psi_I(t)} 
        &= \HH_{\text{eff}}\ket{\psi_I} \\ 
        &= -\hbar \sum_k \ab(g_k^* c_2(t) e^{i(\nu_k - \omega_{12})t} + g_k ^* c_3(t) e^{i(\nu_k - \omega_{13})t})\ket{k} \\
        &\qquad -\hbar\sum_k g_k c_k(t) e^{-i(\nu_k - \omega_{12})t}\ket{2} + g_k c_k(t) e^{-i(\nu_k - \omega_{13})t}\ket{3}
    \end{align*}
    We equate this to $i\hbar(\dot c_2(t)\ket 2 + \dot c_3(t)\ket 3 + \sum_k \dot c_k(t) \ket{k})$ and use the orthogonality of the basis kets to get the system 
    \begin{align*}
        \dot c_2(t) &= i\sum_k g_k c_k(t) e^{-i(\nu_k - \omega_{12})t} \\
        \dot c_3(t) &= i\sum_k g_k c_k(t) e^{-i(\nu_k - \omega_{13})t} \\
        \dot c_k(t) &= i g_k^* c_2(t) e^{i(\nu_k - \omega_{12})t} + ig_k ^* c_3(t) e^{i(\nu_k - \omega_{13})t}
    \end{align*}
\end{ans}
\end{document}